% GENERATED FILE. Edit canonical lessons, then run npm run generate:books.
% canonical-source-hash: fnv1a64:b15c54be4ee163bc
% canonical-lessons: RU-C205-zabyvat, RU-C205-uchitsya, RU-C205-izuchat, RU-C205-prepodavat, RU-C205-sdavat

\chapter{To forget, To be a student, To study, To teach, To hand in}
\label{ch:ru-to-forget-to-be-a-student-to-study-to-teach-to-hand-in}
\hlchaptermodality{205}

\begin{chapteropening}
\textbf{By the end of this chapter:} \emph{I can say to forget, to be a student, to study, to teach and to hand in.}

Complete the last lesson of chapter 205: I can say to forget, to be a student, to study, to teach and to hand in.
\end{chapteropening}

\section[\texorpdfstring{\ru{забывать} (zabyvát)}{забывать (zabyvát)}]{\ru{забывать} (zabyvát) — to forget}

\label{lesson:RU-C205-zabyvat}



\begin{quote}
Before the new one: say the Russian for to cost, then the Russian for to spend.
\end{quote}

\subsection*{You'll want to know: \ru{забывать}}
\textbf{\ru{забывать}} (\emph{zabyvát}) — \textquotedblleft{}to forget\textquotedblright{}.

\begin{cousinweb}
\textbf{\ru{забывать}} is \emph{\ru{за}-}, \textquotedblleft{}beyond\textquotedblright{}, with \emph{-\ru{бывать}}, from \emph{\ru{быть}}, \textquotedblleft{}to be\textquotedblright{} (Vasmer): to forget is to let a thing go past being.
\end{cousinweb}

\begin{grammarlens}[title={What you've built}]
One more word for this chapter.
\end{grammarlens}

\subsection*{Your turn}
\begin{itemize}
\raggedright
  \item \emph{Say it:} \emph{zabyvát}
  \item \emph{Say it:} \emph{zabyvát}, once more
  \item \emph{Say it:} say \emph{zabyvát}
  \item \emph{Recall:} say the Russian for to cost, then the Russian for to spend, then say \emph{zabyvát} again
\end{itemize}

\begin{tcolorbox}[breakable,colback=teal!4,colframe=teal!35!black,title={Before you move on}]
What does \ru{забывать} mean? (\textquotedblleft{}To forget\textquotedblright{}.) Say it once more.
\end{tcolorbox}
\section[\texorpdfstring{\ru{учиться} (uchítsya)}{учиться (uchítsya)}]{\ru{учиться} (uchítsya) — to be a student, to study}

\label{lesson:RU-C205-uchitsya}



\begin{quote}
Before the new one: say the Russian for to spend, then the Russian for to forget.
\end{quote}

\subsection*{You'll want to know: \ru{учиться}}
\textbf{\ru{учиться}} (\emph{uchítsya}) — \textquotedblleft{}to be a student, to study\textquotedblright{}.

\begin{cousinweb}
\textbf{\ru{учиться}} is \emph{\ru{учить}}, \textquotedblleft{}to teach\textquotedblright{}, with \emph{-\ru{ся}}, \textquotedblleft{}oneself\textquotedblright{}. \emph{\ru{учить}} $\leftarrow$ Proto-Slavic \emph{*učiti}, first \textquotedblleft{}to accustom\textquotedblright{} (Vasmer).
\end{cousinweb}

\begin{grammarlens}[title={What you've built}]
One more word for this chapter.
\end{grammarlens}

\subsection*{Your turn}
\begin{itemize}
\raggedright
  \item \emph{Say it:} \emph{uchítsya}
  \item \emph{Say it:} \emph{uchítsya}, once more
  \item \emph{Say it:} say \emph{uchítsya}
  \item \emph{Recall:} say the Russian for to spend, then the Russian for to forget, then say \emph{uchítsya} again
\end{itemize}

\begin{tcolorbox}[breakable,colback=teal!4,colframe=teal!35!black,title={Before you move on}]
What does \ru{учиться} mean? (\textquotedblleft{}To be a student, to study\textquotedblright{}.) Say it once more.
\end{tcolorbox}
\section[\texorpdfstring{\ru{изучать} (izuchát)}{изучать (izuchát)}]{\ru{изучать} (izuchát) — to study (a subject)}

\label{lesson:RU-C205-izuchat}



\begin{quote}
Before the new one: say the Russian for to forget, then the Russian for to be a student.
\end{quote}

\subsection*{You'll want to know: \ru{изучать}}
\textbf{\ru{изучать}} (\emph{izuchát}) — \textquotedblleft{}to study (a subject)\textquotedblright{}.

\begin{cousinweb}
\textbf{\ru{изучать}} is \emph{\ru{из}-}, \textquotedblleft{}out, thoroughly\textquotedblright{}, with \emph{-\ru{учать}}, from \emph{\ru{учить}}, \textquotedblleft{}to teach, to learn\textquotedblright{}: to study a subject through.
\end{cousinweb}

\begin{grammarlens}[title={What you've built}]
One more word for this chapter.
\end{grammarlens}

\subsection*{Your turn}
\begin{itemize}
\raggedright
  \item \emph{Say it:} \emph{izuchát}
  \item \emph{Say it:} \emph{izuchát}, once more
  \item \emph{Say it:} say \emph{izuchát}
  \item \emph{Recall:} say the Russian for to forget, then the Russian for to be a student, then say \emph{izuchát} again
\end{itemize}

\begin{tcolorbox}[breakable,colback=teal!4,colframe=teal!35!black,title={Before you move on}]
What does \ru{изучать} mean? (\textquotedblleft{}To study (a subject)\textquotedblright{}.) Say it once more.
\end{tcolorbox}
\section[\texorpdfstring{\ru{преподавать} (prepodavát)}{преподавать (prepodavát)}]{\ru{преподавать} (prepodavát) — to teach (a subject)}

\label{lesson:RU-C205-prepodavat}



\begin{quote}
Before the new one: say the Russian for to be a student, then the Russian for to study.
\end{quote}

\subsection*{You'll want to know: \ru{преподавать}}
\textbf{\ru{преподавать}} (\emph{prepodavát}) — \textquotedblleft{}to teach (a subject)\textquotedblright{}.

\begin{cousinweb}
\textbf{\ru{преподавать}} is \emph{\ru{пре}-} and \emph{\ru{по}-} with \emph{\ru{давать}}, \textquotedblleft{}to give\textquotedblright{}: to teach is to hand knowledge on. Vasmer takes \emph{\ru{давать}} from \emph{\ru{дать}}, \textquotedblleft{}to give\textquotedblright{}, kin to Lithuanian \emph{dovanà}, \textquotedblleft{}a gift\textquotedblright{}.
\end{cousinweb}

\begin{grammarlens}[title={What you've built}]
One more word for this chapter.
\end{grammarlens}

\subsection*{Your turn}
\begin{itemize}
\raggedright
  \item \emph{Say it:} \emph{prepodavát}
  \item \emph{Say it:} \emph{prepodavát}, once more
  \item \emph{Say it:} say \emph{prepodavát}
  \item \emph{Recall:} say the Russian for to be a student, then the Russian for to study, then say \emph{prepodavát} again
\end{itemize}

\begin{tcolorbox}[breakable,colback=teal!4,colframe=teal!35!black,title={Before you move on}]
What does \ru{преподавать} mean? (\textquotedblleft{}To teach (a subject)\textquotedblright{}.) Say it once more.
\end{tcolorbox}
\section[\texorpdfstring{\ru{сдавать} (sdavát)}{сдавать (sdavát)}]{\ru{сдавать} (sdavát) — to hand in; to take (an exam)}

\label{lesson:RU-C205-sdavat}



\begin{quote}
Before the new one: say the Russian for to study, then the Russian for to teach.
\end{quote}

\subsection*{You'll want to know: \ru{сдавать}}
\textbf{\ru{сдавать}} (\emph{sdavát}) — \textquotedblleft{}to hand in; to take (an exam)\textquotedblright{}.

\begin{cousinweb}
\textbf{\ru{сдавать}} is \emph{\ru{с}-}, \textquotedblleft{}off, over\textquotedblright{}, with \emph{\ru{давать}}, \textquotedblleft{}to give\textquotedblright{}: to hand in, and so to take an exam. Vasmer takes \emph{\ru{давать}} from \emph{\ru{дать}}, \textquotedblleft{}to give\textquotedblright{}, kin to Lithuanian \emph{dovanà}, \textquotedblleft{}a gift\textquotedblright{}.
\end{cousinweb}

\begin{grammarlens}[title={What you've built}]
One more word for this chapter.
\end{grammarlens}

\subsection*{Your turn}
\begin{itemize}
\raggedright
  \item \emph{Say it:} \emph{sdavát}
  \item \emph{Say it:} \emph{sdavát}, once more
  \item \emph{Say it:} say \emph{sdavát}
  \item \emph{Recall:} say the Russian for to study, then the Russian for to teach, then say \emph{sdavát} again
\end{itemize}

\begin{tcolorbox}[breakable,colback=teal!4,colframe=teal!35!black,title={Before you move on}]
What does \ru{сдавать} mean? (\textquotedblleft{}To hand in; to take (an exam)\textquotedblright{}.) Say it once more.
\end{tcolorbox}
